\documentclass[12pt,a4paper,fontset=none]{ctexart}

\usepackage[a4paper,margin=2.4cm]{geometry}
\usepackage{amsmath}
\usepackage{amssymb}
\usepackage{booktabs}
\usepackage{longtable}
\usepackage{array}
\usepackage{enumitem}
\usepackage{hyperref}
\usepackage{xcolor}
\usepackage{listings}

\setCJKmainfont{Microsoft JhengHei}
\setCJKsansfont{Microsoft JhengHei}
\setCJKmonofont{Microsoft JhengHei}

\hypersetup{
  colorlinks=true,
  linkcolor=blue!55!black,
  urlcolor=blue!55!black
}

\lstdefinestyle{cmd}{
  basicstyle=\ttfamily\small,
  breaklines=true,
  frame=single,
  columns=fullflexible
}

\setlist[itemize]{leftmargin=2em}
\setlist[enumerate]{leftmargin=2em}

\title{\textbf{shogiAI 資料結構專題報告}\\
\large 棋譜資料庫、程式資料模型、API JSON 與搜尋資料結構}
\author{專案：shogiAI}
\date{2026 年 5 月}

\begin{document}
\maketitle

\begin{abstract}
本報告針對 shogiAI 專案中的資料結構與資料庫架構進行完整說明。專案處理的資料從原始 CSA 棋譜開始，經過 Python 解析、MySQL 正規化儲存、NumPy 訓練資料集、PyTorch checkpoint、HTTP JSON API，最後在前端以棋盤與手順列表呈現。報告從資料結構角度分析：原始檔案格式如何被拆解、資料庫如何設計、每張表的功能與關聯、Python dataclass 如何表示棋子與局面、後端如何將資料序列化成 JSON、前端如何保存狀態，以及 AI 搜尋如何使用 transposition table、killer moves、history heuristic 等資料結構。此報告重點不是模型準確率，而是「資料如何被組織、保存、傳輸與操作」。
\end{abstract}

\tableofcontents
\newpage

\section{資料結構在本專案中的角色}

shogiAI 並不是只讀取棋譜並顯示棋盤，而是將不同層級的資料結構串成完整系統：

\[
\text{CSA 文字檔}
\rightarrow
\text{Python 物件}
\rightarrow
\text{MySQL 關聯表}
\rightarrow
\text{NumPy 陣列}
\rightarrow
\text{JSON API}
\rightarrow
\text{DOM 棋盤}
\]

因此本專案的資料結構可分成六類：

\begin{enumerate}
  \item 檔案層資料結構：CSA、USI、SFEN、NPZ、PT。
  \item 資料庫層資料結構：users、players、game\_records、moves、positions。
  \item Python 記憶體資料結構：dataclass、dict、list、cshogi.Board。
  \item 機器學習資料結構：state tensor、move label、value label。
  \item API 傳輸資料結構：JSON object、array、nested dictionary。
  \item 搜尋演算法資料結構：transposition table、SearchContext、SearchResult、killer/history tables。
\end{enumerate}

\section{原始棋譜資料結構}

\subsection{CSA 棋譜}

CSA 是文字格式，每一行代表不同類型資料。常見內容包括：

\begin{longtable}{p{0.25\textwidth}p{0.65\textwidth}}
\toprule
\textbf{CSA 行} & \textbf{意義}\\
\midrule
\endhead
\texttt{N+} & 先手棋手名稱。\\
\texttt{N-} & 後手棋手名稱。\\
\texttt{\$EVENT} & 賽事名稱。\\
\texttt{\$SITE} & 比賽地點。\\
\texttt{\$START\_TIME} & 對局日期時間。\\
\texttt{PI} & 標準初始局面。\\
\texttt{P1} 到 \texttt{P9} & 棋盤每一列的棋子配置。\\
\texttt{+7776FU} & 先手一手走法。\\
\texttt{-3334FU} & 後手一手走法。\\
\texttt{\%TORYO} & 終局結果，例如投了。\\
\bottomrule
\end{longtable}

CSA 的優點是完整、可讀；缺點是不適合直接查詢。因此專案需要將 CSA 轉換成資料庫與物件結構。

\subsection{USI 走法}

USI 是程式內部與前後端溝通常用的走法格式，例如：

\begin{itemize}
  \item \texttt{7g7f}：從 7g 走到 7f。
  \item \texttt{2b3c+}：從 2b 走到 3c 並升變。
  \item \texttt{P*5e}：打一枚步到 5e。
\end{itemize}

USI 比 CSA 更適合程式處理，因為它是 \texttt{cshogi} 可直接解析的走法格式。

\subsection{SFEN 局面}

SFEN 是完整局面的文字表示。專案在資料庫 \texttt{positions.sfen} 中保存每一手後的 SFEN。這樣任何局面都可以用：

\begin{lstlisting}[style=cmd]
board = cshogi.Board(sfen)
\end{lstlisting}

直接重建。

\section{資料庫資料結構}

\subsection{設計目標}

資料庫設計目標包含：

\begin{itemize}
  \item 將半結構化 CSA 文字轉成可查詢資料。
  \item 避免棋手名稱與棋局資訊重複儲存。
  \item 保存每手走法與每手後局面。
  \item 讓前端可快速取得任意手數局面。
  \item 讓數據科學流程可從資料庫匯出訓練樣本。
\end{itemize}

\subsection{五張主要資料表}

資料庫由 \texttt{src/import\_csa\_to\_db.py} 建立，包含五張表。

\subsubsection{\texttt{users}}

\texttt{users} 表保存匯入者或未來上傳者資訊。

\begin{longtable}{p{0.25\textwidth}p{0.25\textwidth}p{0.40\textwidth}}
\toprule
\textbf{欄位} & \textbf{型別} & \textbf{說明}\\
\midrule
\endhead
\texttt{user\_id} & INT PK & 使用者主鍵。\\
\texttt{username} & VARCHAR & 使用者名稱。\\
\texttt{email} & VARCHAR & 電子郵件，可唯一。\\
\texttt{password\_hash} & VARCHAR & 保留登入功能擴充。\\
\texttt{role} & VARCHAR & 使用者角色。\\
\texttt{created\_at} & DATETIME & 建立時間。\\
\bottomrule
\end{longtable}

\subsubsection{\texttt{players}}

\texttt{players} 保存棋手資料。

\begin{longtable}{p{0.25\textwidth}p{0.25\textwidth}p{0.40\textwidth}}
\toprule
\textbf{欄位} & \textbf{型別} & \textbf{說明}\\
\midrule
\endhead
\texttt{player\_id} & INT PK & 棋手主鍵。\\
\texttt{player\_name} & VARCHAR & 棋手姓名。\\
\texttt{rank\_name} & VARCHAR & 段位或稱號。\\
\texttt{created\_at} & DATETIME & 建立時間。\\
\bottomrule
\end{longtable}

將棋手獨立成表可避免每盤棋重複儲存棋手名稱，也方便日後統計某棋手對局。

\subsubsection{\texttt{game\_records}}

\texttt{game\_records} 是棋局主表。

\begin{longtable}{p{0.27\textwidth}p{0.23\textwidth}p{0.40\textwidth}}
\toprule
\textbf{欄位} & \textbf{型別} & \textbf{說明}\\
\midrule
\endhead
\texttt{game\_id} & INT PK & 棋局主鍵。\\
\texttt{uploader\_id} & INT FK & 匯入者。\\
\texttt{black\_player\_id} & INT FK & 先手棋手。\\
\texttt{white\_player\_id} & INT FK & 後手棋手。\\
\texttt{event\_name} & VARCHAR & 賽事名稱。\\
\texttt{site} & VARCHAR & 比賽地點。\\
\texttt{opening} & VARCHAR & 開局或戰型。\\
\texttt{result} & VARCHAR & 終局結果。\\
\texttt{source\_format} & VARCHAR & 來源格式，目前為 CSA。\\
\texttt{original\_file\_name} & VARCHAR & 原始檔案名稱，用於追蹤與去重。\\
\texttt{played\_at} & DATE & 對局日期。\\
\texttt{created\_at} & DATETIME & 建立時間。\\
\bottomrule
\end{longtable}

\subsubsection{\texttt{moves}}

\texttt{moves} 保存每一手棋。

\begin{longtable}{p{0.25\textwidth}p{0.25\textwidth}p{0.40\textwidth}}
\toprule
\textbf{欄位} & \textbf{型別} & \textbf{說明}\\
\midrule
\endhead
\texttt{move\_id} & INT PK & 走法主鍵。\\
\texttt{game\_id} & INT FK & 所屬棋局。\\
\texttt{move\_number} & INT & 第幾手。\\
\texttt{side\_to\_move} & ENUM & 先手或後手。\\
\texttt{original\_move} & VARCHAR & 原始走法。\\
\texttt{usi\_move} & VARCHAR & USI 格式走法。\\
\texttt{move\_time} & VARCHAR & 用時資料。\\
\texttt{comment} & TEXT & 棋譜註解。\\
\bottomrule
\end{longtable}

\subsubsection{\texttt{positions}}

\texttt{positions} 保存每一手後的完整局面。

\begin{longtable}{p{0.25\textwidth}p{0.25\textwidth}p{0.40\textwidth}}
\toprule
\textbf{欄位} & \textbf{型別} & \textbf{說明}\\
\midrule
\endhead
\texttt{position\_id} & INT PK & 局面主鍵。\\
\texttt{game\_id} & INT FK & 所屬棋局。\\
\texttt{move\_number} & INT & 第幾手後的局面，0 為初始局面。\\
\texttt{side\_to\_move} & ENUM & 此局面輪到哪一方下。\\
\texttt{sfen} & TEXT & 完整局面表示。\\
\texttt{sfen\_hash} & CHAR(64) & SFEN 的 SHA-256 hash。\\
\texttt{is\_check} & BOOLEAN & 是否被將軍。\\
\texttt{legal\_moves\_count} & INT & 合法手數量。\\
\bottomrule
\end{longtable}

\subsection{資料表關係}

關係如下：

\begin{itemize}
  \item 一個 user 可以上傳多個 game\_records。
  \item 一個 player 可以出現在多個 game\_records。
  \item 一個 game\_record 對應多個 moves。
  \item 一個 game\_record 對應多個 positions。
\end{itemize}

概念圖：

\[
\text{users}
\xrightarrow{1:N}
\text{game\_records}
\xleftarrow{N:1}
\text{players}
\]

\[
\text{game\_records}
\xrightarrow{1:N}
\text{moves}
\qquad
\text{game\_records}
\xrightarrow{1:N}
\text{positions}
\]

\subsection{為什麼 moves 和 positions 分開}

\texttt{moves} 表表示「動作」，\texttt{positions} 表表示「狀態」。這是資料結構上很重要的分離：

\[
\text{Position}_{t}
\xrightarrow{\text{Move}_{t+1}}
\text{Position}_{t+1}
\]

若只保存 moves，查第 $N$ 手局面必須從初始局面 replay 到第 $N$ 手。若保存 positions，前端可以直接查：

\begin{lstlisting}[style=cmd]
SELECT sfen
FROM positions
WHERE game_id = ? AND move_number = ?
\end{lstlisting}

這是用儲存空間換取查詢速度。

\section{Python 記憶體資料結構}

\subsection{\texttt{cshogi.Board}}

\texttt{cshogi.Board} 是專案最底層的規則資料結構。它保存：

\begin{itemize}
  \item 9x9 棋盤。
  \item 雙方持駒。
  \item 目前手番。
  \item 合法手產生器。
  \item SFEN 轉換。
  \item Zobrist hash。
\end{itemize}

常用操作：

\begin{lstlisting}[style=cmd]
board = cshogi.Board()
move = board.move_from_usi("7g7f")
board.is_legal(move)
board.push(move)
board.sfen()
board.legal_moves
\end{lstlisting}

專案不自行實作將棋規則，而是使用 \texttt{cshogi} 保證規則正確性。

\subsection{API 中的 dataclass}

\texttt{src/csa\_browser\_api.py} 定義多個 dataclass。

\subsubsection{\texttt{Piece}}

\begin{lstlisting}[style=cmd]
Piece(color: str, kind: str)
\end{lstlisting}

表示前端棋盤上的一枚棋子：

\begin{itemize}
  \item \texttt{color}：\texttt{+} 或 \texttt{-}。
  \item \texttt{kind}：CSA 棋子代碼，例如 \texttt{FU}、\texttt{HI}、\texttt{OU}。
\end{itemize}

\subsubsection{\texttt{MoveRecord}}

\begin{lstlisting}[style=cmd]
MoveRecord(
  ply,
  color,
  from_square,
  to_square,
  piece,
  usi_like,
  captured
)
\end{lstlisting}

保存一手棋的前端展示資訊。

\subsubsection{\texttt{Position}}

\begin{lstlisting}[style=cmd]
Position(
  ply,
  turn,
  board,
  hands,
  last_move
)
\end{lstlisting}

保存某一手後的完整局面。

\subsubsection{\texttt{Game}}

\begin{lstlisting}[style=cmd]
Game(
  id,
  name,
  metadata,
  positions,
  moves
)
\end{lstlisting}

本機 CSA 模式會將整盤棋解析成 \texttt{Game}，其中 \texttt{positions} 是局面序列。

\section{前後端 JSON 資料結構}

\subsection{棋局清單}

\texttt{GET /api/games} 回傳：

\begin{lstlisting}[style=cmd]
{
  "games": [
    {
      "id": "12",
      "name": "2024 Event A vs B",
      "moves": 105,
      "black": "A",
      "white": "B",
      "event": "...",
      "opening": "...",
      "result": "TORYO",
      "startTime": "2024-01-01"
    }
  ]
}
\end{lstlisting}

\subsection{局面資料}

\texttt{GET /api/games/<id>?ply=N} 回傳：

\begin{lstlisting}[style=cmd]
{
  "game": {...},
  "ply": 30,
  "maxPly": 105,
  "turn": "+",
  "turnLabel": "先手",
  "board": [[...], ...],
  "hands": {
    "+": {"FU": 1, "KA": 0},
    "-": {"FU": 0, "KA": 1}
  },
  "handOrder": ["HI", "KA", "KI", "GI", "KE", "KY", "FU"],
  "lastMove": {...},
  "moves": [...],
  "isCheck": false,
  "legalMovesCount": 42
}
\end{lstlisting}

其中 \texttt{board} 是 9x9 陣列，每格包含：

\begin{lstlisting}[style=cmd]
{
  "square": "77",
  "piece": {
    "color": "+",
    "kind": "FU",
    "label": "步",
    "promoted": false
  }
}
\end{lstlisting}

\subsection{自行對弈請求}

前端送出：

\begin{lstlisting}[style=cmd]
{
  "moves": ["7g7f", "3c3d", "2g2f"]
}
\end{lstlisting}

後端使用 \texttt{replay\_usi\_moves} 重播整個走法序列，再回傳目前局面、合法手與結果。

\subsection{AI 對弈請求}

\begin{lstlisting}[style=cmd]
{
  "moves": ["7g7f", "3c3d"],
  "playerSide": "+",
  "depth": 3,
  "timeLimitMs": 1000
}
\end{lstlisting}

後端回傳 AI 搜尋資訊：

\begin{lstlisting}[style=cmd]
{
  "search": {
    "score": 120,
    "depth": 3,
    "nodes": 5342,
    "timedOut": false,
    "pv": ["2g2f", "8c8d"]
  }
}
\end{lstlisting}

\section{機器學習資料結構}

\subsection{State Tensor}

模型輸入：

\[
states \in \mathbb{R}^{N \times 43 \times 9 \times 9}
\]

這是四維陣列：

\begin{itemize}
  \item $N$：樣本數。
  \item 43：特徵平面。
  \item 9：棋盤列。
  \item 9：棋盤行。
\end{itemize}

\subsection{Move Label}

\[
moves \in \mathbb{Z}^{N}
\]

每個值是 0 到 13688 之間的整數，代表 policy 分類標籤。

\subsection{Value Label}

\[
values \in \mathbb{R}^{N}
\]

每個值是：

\[
-1,\quad 0,\quad 1
\]

\subsection{Metadata}

\texttt{meta} 保存 JSON 字串，用於追蹤樣本來源：

\begin{lstlisting}[style=cmd]
{
  "game_id": 12,
  "ply": 31,
  "sfen": "...",
  "move_usi": "7g7f",
  "result": "TORYO",
  "winner": "black"
}
\end{lstlisting}

這讓模型資料不是黑盒資料，而是可追蹤回原始棋局與手數。

\section{搜尋演算法資料結構}

\subsection{\texttt{SearchResult}}

\begin{lstlisting}[style=cmd]
SearchResult(
  move,
  score,
  depth,
  nodes,
  pv,
  timed_out
)
\end{lstlisting}

保存一次 AI 搜尋的結果：

\begin{itemize}
  \item 最佳手。
  \item 分數。
  \item 完成深度。
  \item 搜尋節點數。
  \item principal variation。
  \item 是否逾時。
\end{itemize}

\subsection{\texttt{SearchContext}}

\texttt{SearchContext} 保存搜尋過程共享狀態：

\begin{itemize}
  \item \texttt{deadline}：時間限制。
  \item \texttt{move\_orderer}：走法排序器，可接 policy model。
  \item \texttt{evaluator}：局面評估函式。
  \item \texttt{nodes}：已搜尋節點數。
  \item \texttt{tt}：transposition table。
  \item \texttt{killers}：killer move table。
  \item \texttt{history}：history heuristic table。
\end{itemize}

\subsection{Transposition Table}

資料結構：

\begin{lstlisting}[style=cmd]
dict[tuple[int, int], TTEntry]
\end{lstlisting}

key 包含：

\begin{itemize}
  \item board zobrist hash。
  \item repetition count。
\end{itemize}

value 是：

\begin{lstlisting}[style=cmd]
TTEntry(depth, score, flag, best_move)
\end{lstlisting}

用途是避免重複搜尋同一局面。

\subsection{Killer Moves}

資料結構：

\begin{lstlisting}[style=cmd]
dict[int, list[int]]
\end{lstlisting}

key 是搜尋深度 ply，value 是該深度曾造成 beta cutoff 的安靜手。搜尋同一層其他分支時，這些手會被優先排序。

\subsection{History Heuristic}

資料結構：

\begin{lstlisting}[style=cmd]
dict[int, int]
\end{lstlisting}

key 是 move，value 是該 move 過去造成剪枝的累積分數。這是一種全域走法排序記憶。

\section{前端狀態資料結構}

前端 \texttt{web/app.js} 使用多組狀態變數。

\subsection{棋譜瀏覽狀態}

\begin{lstlisting}[style=cmd]
currentGameId
currentPly
maxPly
latestState
\end{lstlisting}

用途：

\begin{itemize}
  \item 記錄目前棋局。
  \item 記錄目前播放到第幾手。
  \item 保存後端回傳的最新局面 JSON。
\end{itemize}

\subsection{自行對弈狀態}

\begin{lstlisting}[style=cmd]
selfMoves
selfRedoMoves
selfState
selectedSelfSource
\end{lstlisting}

用途：

\begin{itemize}
  \item \texttt{selfMoves}：已下走法序列。
  \item \texttt{selfRedoMoves}：復原後可重做的走法。
  \item \texttt{selfState}：目前後端回傳局面。
  \item \texttt{selectedSelfSource}：目前選中的棋盤棋子或持駒。
\end{itemize}

\subsection{AI 對弈狀態}

\begin{lstlisting}[style=cmd]
aiMoves
aiState
aiSelectedSource
aiThinking
aiResignedSide
\end{lstlisting}

用途：

\begin{itemize}
  \item 保存 AI 對弈走法序列。
  \item 控制 AI 思考中不可重複點擊。
  \item 保存玩家是否投降。
  \item 管理玩家點棋盤或持駒時的選取狀態。
\end{itemize}

\section{資料結構設計評估}

\subsection{優點}

\begin{enumerate}
  \item \textbf{狀態與動作分離：} positions 保存狀態，moves 保存動作，符合棋局轉移模型。
  \item \textbf{資料可追蹤：} 訓練樣本 metadata 可追蹤回 game id、ply、SFEN 與走法。
  \item \textbf{前後端解耦：} 前端只依賴 JSON 格式，不依賴 MySQL 或 CSA 內部細節。
  \item \textbf{資料來源抽象：} \texttt{GameSource} 使 MySQL 與本機 CSA 檔可替換。
  \item \textbf{搜尋資料結構完整：} TT、killer、history 都能提升 alpha-beta 搜尋效率。
\end{enumerate}

\subsection{限制}

\begin{enumerate}
  \item positions 表會增加儲存空間，因為每手都保存 SFEN。
  \item 前端狀態目前以全域變數管理，功能再擴充時可考慮更結構化狀態管理。
  \item MySQL 尚未在程式中自動建立索引，資料量變大後需要索引優化。
  \item JSON 傳輸每次回傳完整 moves 列表，超長棋局可再做分頁或快取。
\end{enumerate}

\section{建議索引設計}

若資料量增加，建議加入：

\begin{lstlisting}[style=cmd]
CREATE INDEX idx_game_records_played_at
ON game_records(played_at);

CREATE INDEX idx_game_records_original_file_name
ON game_records(original_file_name);

CREATE INDEX idx_moves_game_number
ON moves(game_id, move_number);

CREATE INDEX idx_positions_game_number
ON positions(game_id, move_number);

CREATE INDEX idx_positions_sfen_hash
ON positions(sfen_hash);
\end{lstlisting}

其中 \texttt{idx\_positions\_game\_number} 對前端查第 N 手局面最重要。

\section{結論}

shogiAI 的資料結構設計核心是將棋局拆成「棋局主資料、棋手、走法、局面、模型樣本、API JSON、前端狀態、搜尋快取」等不同層級。這種設計讓系統可以同時滿足資料庫查詢、棋譜瀏覽、模型訓練與 AI 搜尋需求。從資料結構角度看，本專案最重要的設計是 \texttt{moves} 與 \texttt{positions} 分離，因為它把棋類問題明確表示成狀態轉移模型：

\[
S_t + A_{t+1} \rightarrow S_{t+1}
\]

這個模型不只支援前端播放，也支援機器學習中的「局面到下一手」監督式學習。整體而言，shogiAI 是一個以資料結構為基礎，向資料庫、AI 模型與網頁應用延伸的完整系統。

\appendix
\section{附錄：主要資料流}

\begin{lstlisting}[style=cmd]
CSA file
  -> CSA.Parser
  -> cshogi.Board replay
  -> MySQL game_records / moves / positions
  -> export_policy_dataset.py
  -> NumPy states / moves / values
  -> PyTorch model
  -> csa_browser_api.py JSON
  -> web/app.js DOM render
\end{lstlisting}

\end{document}
